\documentclass[11pt, letterpaper]{article}
\usepackage[a4paper, top=2.5cm, bottom=2.5cm, left=2cm, right=2cm]{geometry}
\usepackage[utf8]{inputenc} 
\usepackage[T1]{fontenc} 
\usepackage{fontspec}
\usepackage[greek,english, main = english]{babel} 
\usepackage{fontspec}
\usepackage{booktabs}
\usepackage{graphicx}
\usepackage{hyperref} 
\setlength{\parskip}{0pt}
\setlength{\parindent}{15pt}
\setmainfont{Libertinus Serif}

\begin{document}
\begin{titlepage}
\title{\textbf{\Huge 1η Εργασία στα Νευρωνικά Δίκτυα Βαθιά Μάθηση} \vspace{7 cm}}

\author{\LARGE Φλοριάν Μπυλυκμπάσι \\ \LARGE Τμήμα Ηλεκτρολόγων Μηχανικών και Μηχανικών Υπολογιστών \\ \LARGE ΑΕΜ:11012 \vspace{4 cm}}
\maketitle
\thispagestyle{empty}
\end{titlepage}


\section*{Εισαγωγή - Γενικές Πληροφορίες}
Για την ενδιάμεση εργασία χρησιμοποιήσαμε Visual Studio Code όπου μας ζητούταν η χρήση δύο machine learning αλγορίθμων 
του K-NearestNeighbor και του NearestCentroid για την κατηγοριοποίηση ενος dataset της επιλογής μας. Για να επιτευχθεί αυτό ήταν αναγκαία
η δημιουργία ενός περιβάλλοντος .venv ώστε να μπορούν να εγκατασταθούν οι κατάλληλες βιβλιοθήκες μέσω pip install. \\ \\
Για την κύρια εργασία χρησιμοποιήσαμε Google Colab επειδή ήταν αναγκαία η χρήση της T4 GPU που δίνει η Google καθώς τα μοντέλα μας 
θέλανε αρκετή ώρα για να εκπαιδευτούν. Για την εργασία αυτή έχουν υλοποιηθεί σχεδόν 10 διαφορετικά μοντέλα στα οποία καθε φορά αλλάζουμε κάτι ώστε 
να δούμε πως θα επηρεαστεί το δίκτυο μας (ποσο σωστά θα εκπαιδευτεί, τι ακρίβεια θα έχει). Στο τέλος υπάρχουν τρία παραδείγματα κατηγοριοποίησης.\\ \\
Οι βιβλιοθήκες που έχουν χρησιμοποιηθεί  είναι οι tensorflow, numpy , time , sklearn και matplotlib.pyplot. \\ \\
Από τα δοσμένα dataset επιλέχτηκε η Cifar10 η οποία είναι μια συλλογή από 60.000 32x32 εικόνες
χωρισμένες σε 10 διαφορετικές κλάσεις. Οι κλάσεις αυτές είναι: \\ \\
\begin{enumerate}
    \item Airplane
    \item Automobile
    \item Bird
    \item Cat
    \item Deer
    \item Dog
    \item Frog
    \item Horse
    \item Ship
    \item Truck \\ \\
\end{enumerate} 
Τέλος, είναι χωρισμένη σε ένα training set απο 50.000 εικόνες και σε ένα testing set από 10.000 εικόνες. Για τους σκοπούς της εργασίας χρησιμοποιούμε
ένα \textbf{validation set} 10.000 εικόνων και μόνο 40.000 εικόνες από το training set. \\  \\
Παρακάτω ξεκινάει η ανάλυση πρώτα των δυο κατηγοριοποιητών και έπειτα ανάλυση του κάθε μοντέλου ξεχωριστά. Στο τέλος θα υπάρχουν σχόλια και συμπεράσματα 
για τα αποτελέσματα μας.

\newpage

\section*{1-NearestNeighbor}
Ο κατηγοριοποιητής αυτός για Κ=1 πιάνει μια είκονα την οποία θέλει να ταξινομήσει. Υπολογίζει λοιπόν από αυτήν την ευκλείδια απόσταση
της από όλες τις εικόνες στο training set μας. Από αυτές τις αποστάσεις μετά κρατάει τη μία που είναι πιο κοντά στην δικιά μας εικόνα. 
Έτσι μετά η εικόνα μας παίρνει την κλάση της πιο κοντινής εικόνας (π.χ. αν η πιο κοντινή εικόνα είχε κλάση 'cat' τότε και η δικιά μας θα πάρει την ίδια).\\ \\
Στον κώδικα μας στην αρχή κάνουμε load το dataset μας μέσω της tensorflow. Έπειτα έχουμε την επιλογή είτε να χρησιμοποιήσουμε όλο το dataset μας 
είτε μόνο ένα κομμάτι αυτού. Στο παραδειγμά μας χρησιμοποιούμε όλο το dataset αλλά ο χρήστης έχει την επιλογή να μειώσει τα δείγματα ώστε να δει τις αλλαγές
στην ακρίβεια. Επειδή ο αλγόριθμος δεν δέχεται 2D εικόνες κάνουμε flatten τις εικόνες σε ένα μεγάλο διάνυσμα 3072 στοιχειών (32x32x3).
Έπειτα μέσω των βοηθητικών εντολών των βιβλιοθήκων μας γίνεται η εκπαίδευση και η αξιολόγηση και στο τελός εκτυπώνονται τα αποτελέσματα. \\ \\ \\ \\


\section*{3-NearestNeighbor}
Παρόμοια με πριν ο κατηγοριοποιητής πιάνει μία εικόνα και υπολογίζει την ευκλείδια απόσταση της από όλες τις υπόλοιπες. Τώρα η διαφορά είναι
ότι κρατάει τις 3 πιο κοντινές αποστάσεις και κοιτάζει τις κλάσεις αυτών των εικόνων. Στο τελός θα επιλέξει την πλειοψηφία αυτών των κλάσεων (π.χ. αν έχω 2 κλάσεις 'cat' 
και 1 'frog', επιλέγει 'cat'.) \\ \\
Ο κώδικας για Κ=3 βρίσκεται στον ίδιο κώδικα με αυτόν για Κ=1  απλώς προσθέσαμε την εκπαίδευση και την αξιολόγηση στις βοηθητικές εντολές της βιβλιοθήκης
για Κ=3 τώρα και στο τέλος τυπώνουμε τα αποτλέμσατα. \\ \\ \\ \\ 

\section*{Αποτελέσματα για NearestNeighbor}
\vspace{3 em}
\begin{figure}[h]  
    \centering
    \includegraphics[width=0.4\textwidth]{knn1.png}
    \caption{Αποτελέσματα για Κ=1}
\end{figure}

\begin{figure}[h]  
    \centering
    \includegraphics[width=0.39\textwidth]{knn3.png}
    \caption{Αποτελέσματα για Κ=3}
\end{figure}

\newpage

Παρατηρούμε ότι και τα δύο έχουν παρόμοιους χρόυνος εκπαίδευσης , γύρω στα 90 δευτερόλεπτα. Ωστόσο για Κ=1 έχουμε κατα 2\% μεγαλύτερη ακρίβεια από ότι για Κ=3.
Οπότε μπορούμε εύκολα να συμπεράνουμε ότι για ολόκληρο το dataset  για Κ=1 ο αλγόριθμος δουλεύει καλύτερα από ότι για Κ=3. \\ \\ \\ \\ 

\section*{Nearest Centroid}
Η βασική ιδέα τώρα εδώ είναι ότι κάθε κλάση πλέον αντιπροσωπεύεται από το κέντρο των δειγμάτων της στον χώρο των χαρακτηριστικών.
Στην αρχή λοιπόν ο αλγόριθμος υπολογιίζει το μέσο όρο όλων των χαρακτηριστικών των δειγμάτων. Όταν λοιπόν έρθει ένα νέο δείγμα 
τότε υπολογίζουμε την απόσταση του από όλα αυτά τα κέντρα που έχουμε. Το δείγμα μετά θα πάρει την κλάση της μικρότερης απόστασης. \\ \\

Ο κώδικας χρησιμοποιεί ακριβώς την ίδια λόγική με πριν , το μόνο που αλλάζει πλέον είναι οτι χρησιμοποιούνται οι βοηθητικές εντολές
της βιβλιοθήκης για την NearestCentroid από οτι την NearestNeighbor για την εκπαίδευση και την αξιολόγηση. \\ \\

\section*{Αποτελέσματα για NearestCentroid}

\vspace{3 em}
\begin{figure}[h]  
    \centering
    \includegraphics[width=0.39\textwidth]{cn.png}
    \caption{Αποτελέσματα για NearestCentroid}
\end{figure}

Βλέπουμε εδώ ότι ο χρόνος εκπαίδευσης είναι εξαιρετικά μικρός ωστόσο η ακρίβεια είναι αρκετά πιο χαμηλή σε σχέση με τον NearestNeighbor (7\% πιο κάτω). \\ \\

\section*{Συμπεράσματα-Σχόλια}
Μπορούμε εύκολα να συμπεράνουμε  αντικειμενικά ότι ο καλύτερος απο τους 3 είναι ο 1-NearestNeighbor καθώς είναι πιο αποδοτικός και ας έχει ελάχιστα μεγαλύτερο χρόνο απο τον 3-NearestNeighbor.
Ο πιο γρήγορος είναι ο NearestCentroid , ωστόσο λόγω της κακής του απόδοσης καθίσταται τελευταίο της λίστας. Ο 3-NearestNeighbor λόγω της χαμηλότερης κατά 2\% απόδοση , βρίσκεται κάτω απο τον 1-
NearestNeighbor.

\newpage

\section*{Υλοποίηση Νευρωνικού Δικτύου για την Cifar10}
Στο σημείο αυτό ακολουθεί η κύρια ιδέα της εργασίας που είναι η υλοποιήση ενός νευρωνικού δικτύου εμπρόσθιας τροφοδότησης το οποίο θα μπορούσε να ήταν είτε MLP (πλήρως συνδεδεμένο)
είτε CNN (συνελικτικό) είτε συνδυασμός των δύο. Παρακάτω υπάρχουν υλοποιήσεις 2 μοντέλων MLP και αρκετών μοντέλων συνδυασμών. Πριν από αυτά όμως θα αναφερθούν κάποιες γενικές πληροφορίες για το 
κάθε νευρωνικό δίκτυο και το πως δουλεύει το καθένα. Επίσης θα σχολιαστούν και θα αναλυθούν κάποιοι σημαντικοί όροι οι οποίοι είναι αναγκαίοι για την κατανόηση των μοντέλων μας.  \\ \\ \\ 
\section*{Multi Layer Perceptrons}
Όταν λέμε εμπρόσθια τροφοδότηση εννοούμε ότι η πληροφορία κινείται προς την μια κατεύθυνση από το επίπεδο εισόδου , στα όσα κρυφά επίπεδα υπάρχουν και τέλος στο επίπεδο εξόδου χωρίς να 
υπάρχει κάποια επανάληψη ή ανάδραση. \\ \\
Τα MLP είναι ένα είδος ενός νευρωνικού δικτύου εμπρόσθιας τροφοδότησης τα οποία αποτελούνται από πλήρως συνδεδεμένα πυκνά επίπεδα τα οποία μετατρέπουν δεδομένα από μια διάσταση σε μια άλλη.
Λέγονται πολύ-επίπεδα επειδή περιέχουν ένα επίπεδο εισόδου ,κανένα, ένα, ή πολλά κρυφά επίπεδα και ένα επίπεδο εξόδου. \\ \\
Στο κρυφό επίπεδο ο καθε νευρώνας υπολογίζει το άθροισμα των βαρών της εισόδου συν ακόμα ένα bias όρο. Έπειτα αυτό το άθροισμα περνάει μέσα από μία συνάρτηση ενεργοποίησης ώστε να εισάγει την μη 
γραμμικότητα που θέλουμε.\textbf{Στα παραδείγματα μας χρησιμοποιούμε την ReLU ως συνάρτηση ενεργοποίησης.}
Αφού το δικτύο παράξει μια έξοδο τότε ακολουθεί ο υπολογισμός του σφάλματος μέσω της χρήσης κάποιας συνάρτησης σφάλματος.
Στα παραδείγματα μας  χρησιμοποιείται η Sparse Categorical Crossentropy. \\ \\
Μετά απο αυτά ακολουθεί ο backpropagation αλγόριθμος ο οποίος αυτό που κάνει είναι να ελαχιστοποιεί την διαφορά μεταξύ της προβλεπόμενης και της πραγματικής εξόδου. Ο τρόπος που δουλεύει
είναι να παίρνει τα σφάλματα προς τα πίσω στο δίκτυο, χρησιμοποιώντας τον κανόνα αλυσίδας του λογισμού που υπολογίζει τα ανάδελτα, και μετά αναδρομικά ενημερώνει τα βάρη και τα biases. \\ \\
Τέλος αυτά τα δίκτυα βασίζονται πάρα πολύ σε αλγορίθμους βελτιστοποίησης τα οποία αναδρομικά βελτιώνουν τα βάρη και τα biases την ώρα της εκπαίδευσης.

\newpage

\section*{Convolutional Neural Networks}
Τα συνελικτικά νευρωνικά δίκτυα έχουν έναν λίγο πιο περίπλοκο σχεδιασμό. Αποτελούνται και αυτά πάλι από ένα επίπεδο εισόδου αλλά τώρα το κρυφό επίπεδο
αποτελείται από ένα συνελικτικό επίπεδο, από ένα pooling επίπεδο και ένα πλήρως συνδεδεμένο επίπεδο (μέσα σε αυτά υπάρχει και το επίπεδο ενεργοποίησης). Στο τέλος κλασικά υπάρχει το γνωστό όπως πριν 
επίπεδο εξόδου. \\ \\
Στην αρχή το επίπεδο είσοδου είναι αυτό που όπως πριν δίνουμε την είσοδο στο μοντέλο μας το οποίο κρατάει τις διαστάσεις όπως είναι. \\ \\
Έπειτα ακολουθεί πρώτα το συνελικτικό επίπεδο στο οποίο προσπαθούμε να εξάγουμε την μορφή της εικόνας. Σε αυτό το επίπεδο εφαρμόζονται 
ενα σετ απο φίλτρα στις εικόνες μας τα οποία συνήθως είναι μικροί πίνακες. Αυτό που εξάγεται από το επίπεδο αυτό είναι οι διαστάσεις που είχαμε πριν
με μονή αλλαγή την τελευταία διάσταση στην οποία τώρα θα μπει ο αριθμός των φίλτρων. \\ \\
Μετά έχουμε ένα επίπεδο ενεργοποίησης στο οποίο παρόμοια με τα MLP εισάγουμε την συνάρτηση ενεργοποίησης στην έξοδο του προηγούμενου επιπέδου.
\textbf{Και εδώ γίνεται χρήση της ReLU στα παραδείγματα μας.} \\ \\
Ύστερα  έχουμε το pooling επίπεδο το οποίο η κύρια λειτουργία του είναι να μειώσει τις δύο πρώτες διαστάσεις μας κάνοντας έτσι 
τους υπολογισμούς γρηγορότερος ενώ ταυτόχρονα μειώνει την μνήμη και μπορούμε να αποφύγουμε την υπερεκπαίδευση. Εμείς κάνουμε χρήση του max pooling. \\ \\
Πριν περάσει η πληροφορία στο πλήρως συνδεδεμένο επίπεδο κάνουμε flatten την πληροφορία σε 1D διάνυσμα καθώς είναι ο μόνος τρόπος ώστε να 
εισέλθει στο επίπεδο αυτό. Εκεί κάνει τον τελευταίο απαραίτητο υπολογισμό κατηγοριοποίησης πριν περάσει στο επίπεδο εξόδου. \\ \\
Στο επίπεδο εξόδου χρησιμοποιούμε την συνάρτηση softmax η οποία μετατρέπει την έξοδο αυτή που πήραμε στο πιθανοτικό αποτέλεσμα κάθε κλάσης. \\ \\
Οι αλγόριθμοι βελτιστοποίησης καθώς και ο αλγόριθμος backpropagation έχουν ακριβώς την ίδια λογική με αυτήν στα MLP. \\ \\
Παρακάτω ακολουθεί η ανάλυση βασικών ορών καθώς είναι πράγματα τα οποία χρησιμοποιούνται αρκετά στα μοντέλα μας.

\newpage

\section*{Βασικοί Όροι}
\begin{itemize}
    \item Overfit(Υπερεκπαίδευση): Συμβαίνει όταν το μοντέλο μαθαίνει πάρα πολλά από τα δεδομένα εκπαίδευσης. Έτσι δούλευει καλά το μοντέλο σε αυτά τα δεδομένα όχι όμως όταν δοκιμάζεται σε νέα δεδομένα. Συνήθως συμβαίνει όταν το μοντέλο είναι πολύπλοκο ή λόγω του μεγέθους των δεδομένων εκπαίδευσης.
    \item Underfit(Υποεκπαίδευση): Εδώ έχουμε το αντίθετο του overfit, διότι το μοντέλο δεν μπορεί να καταλάβει τι γίνεται με τα δεδομένα. Έτσι το μοντέλο δεν δουλεύει καλά ούτε στην εκπαίδευση ούτε στην δοκιμή. Λόγοι μπορεί να είναι η απλότητα του μοντέλου, ανεπαρκές μέγεθος των δεδομένων εκπαίδευσης, ή υπερβολική χρήση τεχνικών κανονικοποίησης που βοηθούν στην αποφυγή υπερκπαίδευσης.
    \item Epoch(Εποχές): Είναι το πόσες φορές το μοντέλο σου θα επεργάζεται τα δεδομένα σου. Στα παραδείγματα κυρίως χρησιμοποιούμε 100.
    \item Batch size: Πόσα δείγματα θα επεξεργαστούν πριν ενημερωθεί το μοντέλο μας. Στα παραδείγματα χρησιμοποιούμε κυρίως 64.
    \item Validation data(Δεδομένα Επικύρωσης): Είναι ένα κομμάτι απο δεδομένα διαφορετικά απο το βασικο training set τα οποία χρησιμοποιούνται ώστε να δούμε πόσο καλά γενικεύεται το μοντέλο σε νέα δεδομένα ώστε να ρυθμίσουμε κατάλληλα τις παραμέτρους εκπαίδευσης.
    \item Early Stopping(Πρόωρη Διακοπή): Είναι μια τεχνική η οποία σταματά την εκπαίδευση πριν ολοκληρωθούν ολα τα epochs επειδή το μοντέλο δεν βελτιώνεται άλλο ή αρχίζει να χειροτερεύει λόγω overfit.
    \item Learning Rate(Ρυθμός Μάθησης): Είναι το πόσο μεγάλα βήματα κάνει το δίκτυο μας όταν ενημερώνει τα βάρη κατα την εκπαίδευση. Όταν είναι μεγάλος ο ρυθμός το μοντέλο μαθαίνει γρήγορα όμως μπορεί να γίνει ασταθές ενώ όταν είναι μικρό το μοντέλο μαθαίνει πιο αργά αλλά σταθερά.
    \item Dropout: Είναι τεχνική που απενεργοποιεί τυχαία ένα ποσοστό νευρώνων βοηθώντας έτσι να αποτρέψουμε την υπερεκπαίδευση. Το μοντέλο δηλαδή  εκπαιδεύεται απο συγκεκριμένους νευρώνες.
    \item Batch Normalization: Τεχνική που ομαλοποιεί(κανονικοποιεί) τις τιμές των νευρώνων μέσα σε κάθε batch κατά την εκπαίδευση. Με αυτήν την τεχνική σταθεροποιείται η εκπαίδευση, επιτρέπεται η χρήση μεγαλύτερου learning rate και αυξάνεται η ταχύτητα εκπαίδευσης.
    \item Weight decay ή L2 Regularization: Είναι και αυτή τεχνική που πρσθέτει μια μικρή ποινή στα μεγάλα βάρη του μοντέλου ώστε να μην μεγαλώσουν υπερβολικά. Στόχος είναι η μείωση υπερκπαίδευσης και η σταθεροποίηση του μοντέλου. \\ \\
\end{itemize}
Αυτοί ήταν κυρίως οι όροι οι οποίοι χρήζαν ανάλυσης καθώς στη συνέχεια χρησιμοποιούνται συνέχεια στα μοντέλα μας. Επίσης θα βοηθήσει πολύ αυτή η ανάλυση των όρων και στην
ανάλυση του κώδικα που ακολουθεί μετά.

\newpage

\section*{Ανάλυση Κώδικα}
Ξεκινάμε κάνοντας import όλες τις απαραίτητες βιβλιοθήκες που χρησιμοποιούνται μετά στον κώδικα . Πολλά απο αυτά που υπάρχουν στον κώδικα μας έχουν ήδη 
αναφερθεί ήδη πάνω. \\ \\
Κάνουμε load την cifar10 και κανονικοποιούμε απευθείας τα pixel και κάνουμε flatten  τις ετικέτες ώστε να είναι η διαδικασία πιο εύκολη. Έπειτα χωρίζουμε τα δεδομένα μας σε training set, validation set
και test set. Το training set εχει 40.000, το validation έχει 10.000 και το test μένει όπως ήταν με 10.000 εικόνες. \\ \\
Μετά ξεκινάει η κατασκευή του μοντέλου μας, όπου παίρνουμε αρχικά τον αριθμό των κλάσεων μας. Έπειτα προστίθενται τα διάφορα επίπεδα μας αναλόγως με το τι έχουμε στο μυαλό μας να κατασκευάσουμε.
Η λογική της κατασκευής αυτής βασίζεται στις πληροφορίες που δώσαμε πριν για τα δύο είδη νευρωνικών δικτύων που μας δόθηκαν. Γενικά τα παραδείγματα ξεκινάνε από απλά μοντέλα και όσο προχωράμε γίνονται πιο πολύπλοκα.\\ \\
Στην συνέχεια κάνουμε compile το μοντέλο μας χρησιμοποιώντας τον Adam Optimizer και την Sparse Categorical Crossentropy για όλα τα παραδείγματα μας. Επίσης έχουμε χρήση της Early Stopping αν παρατηρήσουμε ότι δεν γίνεται
βελτίωση για 25 εποχές ενώ παράλληλα αποθηκεύουμε τα καλύτερα βάρη. Μετά απο αυτά έχουμε προσθέσει το ReduceLROnPlateau
το οποίο αν δει οτι το μοντέλο χειροτερεύει , χαμηλώνει λίγο το learning rate ώστε να σταθεροποίησει κάπως το μοντέλο μας. \\ \\
Αφού περάσουμε από όλα αυτά φτάνουμε στο σημείο που ξεκινάει η εκπαίδευση του μοντέλου μας με τις εποχές και το batch size που αναφέραμε πριν. Επίσης υπολογίζουμε ΜΟΝΟ μια φορά την ακρίβεια και το σφάλμα
σε ολόκληρο το test set.\\ \\
Στο τέλος τυπώνουμε όλες τις απαραίτητες πληροφορίες που ζητάει η άσκηση και 2 επιπρόσθετα γραφήματα που δείχνουν την σχέση μεταξύ της ακρίβειας εκπαίδευσης με την ακρίβεια επικύρωσης και την
σχέση μεταξύ σφάλματος εκπαίδευσης με το σφάλμα επικύρωσης αντίστοιχα. \\ \\ 
Όσο περνάνε τα παραδείγματα θα δούμε στον κώδικα ότι προστίθενται οι διάφορες τεχνικές κανονικοποίησης που αναλύθηκαν πριν. 
Ας δούμε τώρα αναλυτικά ένα-ένα τα μοντέλα μας.

\newpage

\section*{Παράδειγμα-Μοντέλο 1}
\begin{itemize}
    \item Σε αυτό το παράδειγμα κάνουμε χρήση MLP με δύο κρυφά επίπεδα στα οποία βάζουμε απο 64 νευρώνες στο καθένα.
    \item Δεν κάνουμε ΚΑΜΙΑ χρήση κάποιας τεχνικής κανονικοποίησης ώστε να βοηθήσουμε το μοντέλο μας, μόνο του ReduceLROnPlateau που η λειτουργία του εξηγήθηκε πιο πάνω.
    \item Έχουμε 100 εποχές , 64 batch size και Adam Optimizer. \\ 
\end{itemize}
Ο κώδικας μας τυπώνει:
\begin{itemize}
    \item "Total training time = 212.87 seconds"
    \item "Test accuracy: 0.4802 , Test loss: 1.4668" \\ 
\end{itemize}
και έχουμε τα εξής διαγράμματα:
\begin{figure}[h]  
    \centering
    \includegraphics[width=0.5\textwidth]{1accuracy.png}
    \caption{Διάγραμμα ακρίβειας 1ου παραδείγματος}
\end{figure}
\begin{figure}[h]  
    \centering
    \includegraphics[width=0.5\textwidth]{1loss.png}
    \caption{Διάγραμμα σφάλματος 1ου παραδείγματος}
\end{figure}
\newpage
Βλέπουμε ότι το μοντέλο απο το διάγραμμα της ακρίβειας στις πρώτες εποχές  εκπαιδεύεται ραγδαία που σημαίνει οτι μαθαίνει γρήγορα τα βασικά χαρακτηριστικά και μοτίβα των δεδομένων. 'Επειτα 
οι δυο καμπύλες ανεβαίνουν σταδιάκα σχεδόν μαζί εως τις 40 εποχές. Απο εκει και μέχρι τις 100 εποχές παρατηρείται ένα κενό ανάμεσα στις δύο καμπύλες ενώ μάλιστα σταθεροποιούνται και οι
δύο σε κάποιες τιμές(και στα δυο διαγράμματα μας) , γεγονός που δηλώνει οτι απο εκεί και πέρα το μοντέλο μας κανει overfit. \\ \\
Η ακρίβεια είναι σχετικά μικρή (κάτω απο 50\%) διότι το μοντέλο μας είναι σχετικά απλό και αποτελείται απο σχετικά μικρό αριθμό νευρώνων.

\section*{Παράδειγμα-Μοντέλο 2}
\begin{itemize}
    \item Σε αυτό το παράδειγμα κάνουμε χρήση MLP με δύο κρυφά επίπεδα στα οποία βάζουμε απο 32 νευρώνες στο καθένα αυτή τη φορά.
    \item Δεν κάνουμε πολλές χρήσεις  τεχνικών κανονικοποίησης , μόνο του Dropout(30\%) και του ReduceLROnPlateau .
    \item Έχουμε 100 εποχές , 64 batch size και Adam Optimizer. \\ 
\end{itemize}
Ο κώδικας μας τυπώνει:
\begin{itemize}
    \item "Total training time = 191.44 seconds"
    \item "Test accuracy: 0.3400 , Test loss: 1.7964" \\ 
\end{itemize}
και έχουμε τα εξής διαγράμματα:
\begin{figure}[h]  
    \centering
    \includegraphics[width=0.5\textwidth]{2accuracy.png}
    \caption{Διάγραμμα ακρίβειας 2ου παραδείγματος}
\end{figure}
\newpage
\begin{figure}[h]  
    \centering
    \includegraphics[width=0.5\textwidth]{2loss.png}
    \caption{Διάγραμμα σφάλματος 2ου παραδείγματος}
\end{figure}
Πάλι εδώ βλέπουμε στις 10 πρώτες εποχές μια ραγδαία αύξηση στην ακρίβεια στις καμπύλες μας δείχνοντας οτι το μοντέλο μας μαθαίνει γρήγορα τα βασικά χαρακτηριστικά.
Μετά οι δύο καμπύλες μας  φαινεται ότι ανεβαίνουν με ενα πιο αργό και σταθερό ρυθμό και σε αντίθεση με πριν ανεβαίνουν σχετικά μαζί , δηλαδή οι καμπύλες είναι κοντά η μια στην άλλη.
Όμοια και αντίθετα τα ίδια συμβαίνουν και στις καμπύλες σφάλματος. Δεν έχουμε λοιπόν overfit , αλλά αντίθετα το μοντέλο μαθαίνει πολύ αποδοτικά.
Στο τέλος στις τελευταίες εποχές αρχίζουν και σταθεροποιούνται οι καμπύλες της ακρίβειας γύρω στο 0.34 και του σφάλματος κοντα στο 1.80.\\ \\
Έχουμε πιο μικρή ακρίβεια απο πριν επειδή χρησιμοποιούμε λιγότερους νευρώνες και επίσης η Dropout τεχνική απενεργοποιεί παράλληλα κάποιους απο αυτούς τους νευρώνες μας. \\ \\
Παρακάτω θα δούμε μοντέλα πιο πολύπλοκα που θα αυξήσουν κατα πολύ την ακρίβεια ωστόσο θα προκαλέσουν σημαντικά overfit. Καθώς προχωράμε θα προσθέτουμε layers και θα βρούμε τρόπους
να σταματήσουμε αυτό το πρόβλημα του overfit κρατώντας την ακρίβεια ψηλά.

\section*{Παράδειγμα-Μοντέλο 3}
\begin{itemize}
    \item Σε αυτό το παράδειγμα στο κρυφό επίπεδο κάνουμε χρήση 2 επιπέδων Conv2D (2 συνελικτικά επίπεδα), το καθένα να έχει ένα Βatch Νormalization επίπεδο και μετά από όλα αυτά το MaxPooling επίπεδο. 
    Η έξοδος αυτών βρίσκει ενα Fully Connected Layer 32 νευρώνων
    \item Δεν κάνουμε χρήσεις  τεχνικών κανονικοποίησης , μόνο του ReduceLROnPlateau .
    \item Έχουμε 100 εποχές , 64 batch size και Adam Optimizer. \\ 
\end{itemize}
Ο κώδικας μας τυπώνει:
\begin{itemize}
    \item "Total training time = 114.41 seconds"
    \item "Test accuracy: 0.6525 , Test loss: 1.0197" \\ 
\end{itemize}
και έχουμε τα εξής διαγράμματα:

\newpage

\begin{figure}[h]  
    \centering
    \includegraphics[width=0.5\textwidth]{3accuracy.png}
    \caption{Διάγραμμα ακρίβειας 3ου παραδείγματος}
\end{figure}
\begin{figure}[h]  
    \centering
    \includegraphics[width=0.5\textwidth]{3loss.png}
    \caption{Διάγραμμα σφάλματος 3ου παραδείγματος}
\end{figure}
Ξεκινάμε πάλι στην ακρίβεια με ραγδαία αύξηση απο το 0.4 στο 0.55 περίπου στις πολύ πρώτες εποχές.
Η καμπύλη του training πηγαίνει εξαιρετικά γρήγορα γύρω στο 0.9 ενώ η validation καμπύλη σιγά και σταθερά
στο 0.65. Παρατηρούμε ένα τεράστιο κενό μεταξύ των καμπυλών (είτε ακρίβειας είτε σφάλματος) πράγμα που δηλώνει
σοβαρό overfit. Το μοντέλο παύει να εκπαιδεύεται και κάπου στις 30 εποχές καλείται η Early Stopping.\\ \\
Αυτό που καταφέραμε ωστόσο είναι σε σχέση με πριν είναι να αυξήσουμε την ακρίβεια κατά ενα πολύ σημαντικό ποσοστό
(σχέδον κατά 20\%).

\newpage

\section*{Παράδειγμα-Μοντέλο 4}
\begin{itemize}
    \item Σε αυτό το παράδειγμα στο κρυφό επίπεδο κάνουμε χρήση 4 επιπέδων Conv2D (4 συνελικτικά επίπεδα), το καθένα να έχει ένα Βatch Νormalization επίπεδο και 2 MaxPooling επίπεδα (ένα ανα 2 Conv2D επίπεδα). 
    Η έξοδος αυτών βρίσκει ενα Fully Connected Layer 128 νευρώνων
    \item Δεν κάνουμε χρήσεις  τεχνικών κανονικοποίησης , μόνο του ReduceLROnPlateau .
    \item Έχουμε 100 εποχές , 64 batch size και Adam Optimizer. \\ 
\end{itemize}
Ο κώδικας μας τυπώνει:
\begin{itemize}
    \item "Total training time = 168.45 seconds"
    \item "Test accuracy: 0.7609 , Test loss: 0.8037" \\ 
\end{itemize}
και έχουμε τα εξής διαγράμματα:
\begin{figure}[h]  
    \centering
    \includegraphics[width=0.5\textwidth]{4accuracy.png}
    \caption{Διάγραμμα ακρίβειας 4ου παραδείγματος}
\end{figure}
\begin{figure}[h]  
    \centering
    \includegraphics[width=0.5\textwidth]{4loss.png}
    \caption{Διάγραμμα σφάλματος 4ου παραδείγματος}
\end{figure}

\newpage
Στην αρχή ισχύει το ίδιο με όλα τα προηγούμενα παραδείγματα ωστόσο αυτή τη φορά πηγαίνουμε
απο το 0.55 στο 0.7 περιπου. Η καμπύλη του training αρχίζει και αποκλίνει παρα πολύ απότομα ήδη
απο τις πρώτες εποχές και καταφέρνει και φτάνει μεχρι και 1.0 (πλήρης απομνημόνευση). Η άλλη καμπύλη
σταθεροποιείται κάπου στο 0.77 οπότε έχουμε πάλι αρκετό overfit και για αυτό καλείται η Early Stopping
κάπου στις 30 εποχές.\\ \\
Καταφέραμε ωστόσο με την αύξηση ενός μπλοκ να κερδίσουμε πάλι αρκετά στην ακρίβεια κατα ενα ποσοστό περίπου 
10\%.

\section*{Παράδειγμα-Μοντέλο 5}
\begin{itemize}
    \item Σε αυτό το παράδειγμα στο κρυφό επίπεδο κάνουμε χρήση 6 επιπέδων Conv2D (6 συνελικτικά επίπεδα), το καθένα να έχει ένα Βatch Νormalization επίπεδο και 3 MaxPooling επίπεδα (ένα ανα 2 Conv2D επίπεδα). 
    Η έξοδος αυτών βρίσκει ενα Fully Connected Layer 256 νευρώνων
    \item Δεν κάνουμε χρήσεις  τεχνικών κανονικοποίησης , μόνο του ReduceLROnPlateau .
    \item Έχουμε 100 εποχές , 64 batch size και Adam Optimizer. \\ 
\end{itemize}
Ο κώδικας μας τυπώνει:
\begin{itemize}
    \item "Total training time = 190.41 seconds"
    \item "Test accuracy: 0.7682 , Test loss: 0.6907" \\ 
\end{itemize}
και έχουμε τα εξής διαγράμματα:
\begin{figure}[h]  
    \centering
    \includegraphics[width=0.5\textwidth]{5accuracy.png}
    \caption{Διάγραμμα ακρίβειας 5ου παραδείγματος}
\end{figure}
\newpage
\begin{figure}[h]  
    \centering
    \includegraphics[width=0.5\textwidth]{5loss.png}
    \caption{Διάγραμμα σφάλματος 5ου παραδείγματος}
\end{figure}
Η κατάσταση εδώ είναι ακριβώς πανομοιότυπη με πριν οπότε δεν έχουμε πολλά να εξηγήσουμε. Η ακρίβεια αυξήθηκε κατά ελάχιστα και το σφάλμα 
κατάφερε να πέσει κατα 0.1. Συμπεραίνουμε λοιπόν ότι τα 6 συνελικτικά επίπεδα αυτά (τα 3 μπλοκ μαζί δηλαδή με όλα τα άλλα επίπεδα) είναι το "όριο" για το μοντέλο
και πως δεν χρειάζεται να το κάνουμε πιο πολύπλοκο. Παρακάτω λοιπόν θα εισάγουμε τεχνικές που θα σταματήσουν το overfit κρατώντας και την ακρίβεια ωστόσο ψηλά.

\section*{Παράδειγμα-Μοντέλο 6}
\begin{itemize}
    \item Σε αυτό το παράδειγμα στο κρυφό επίπεδο κάνουμε χρήση 6 επιπέδων Conv2D (6 συνελικτικά επίπεδα), το καθένα να έχει ένα Βatch Νormalization επίπεδο και 3 MaxPooling επίπεδα (ένα ανα 2 Conv2D επίπεδα). 
    Η έξοδος αυτών βρίσκει ενα Fully Connected Layer 256 νευρώνων
    \item Εκτός του ReduceLROnPlateau, πλέον βάζουμε και ένα Dropout(75\%) πριν την έξοδο.
    \item Έχουμε 100 εποχές , 64 batch size και Adam Optimizer. \\ 
\end{itemize}
Ο κώδικας μας τυπώνει:
\begin{itemize}
    \item "Total training time = 233.17 seconds"
    \item "Test accuracy: 0.8100 , Test loss: 0.6199" \\ 
\end{itemize}
και έχουμε τα εξής διαγράμματα:

\newpage

\begin{figure}[h]  
    \centering
    \includegraphics[width=0.5\textwidth]{6accuracy.png}
    \caption{Διάγραμμα ακρίβειας 6ου παραδείγματος}
\end{figure}
\begin{figure}[h]  
    \centering
    \includegraphics[width=0.5\textwidth]{6loss.png}
    \caption{Διάγραμμα σφάλματος 6ου παραδείγματος}
\end{figure}
Τα πράγματα εξακολουθούν να είναι τα ίδια στην αρχή και γενικά ως την 10η εποχή περίπου έχουμε ένα καλό ξεκίνημα με το μοντέλο να εκπαιδεύεται 
σωστά και να φτάνει ως και το 0.78 σχεδόν. Μετά υπάρχει μια σχετικά ελεγχόμενη απόκλιση σε σχέση με πριν περίπου 13\% με την training καμπύλη να αγγίζει
το 0.95 και την validation καμπύλη το 0.82. Ώστοσο μετά υπάρχει μια σταθεροποίηση αναγκάζοντας έτσι την κλήση της Early Stopping ώστε να αποθηκευτούν τα 
καλύτερα βάρη. \\ \\
Το μόνο σίγουρο είναι ότι αυτό το μοντέλο είναι σίγουρα καλύτερο απο τα προηγούμενα καθώς καταφέραμε και αγγίξαμε ακρίβεια πάνω απο 0.8. Ωστόσο
το overfit ακόμα παραμένει που σήμαινει ότι πρέπει κάτι ακόμα να προσθέσουμε ώστε να βοηθήσουμε το μοντέλο μας.

\newpage

\section*{Παράδειγμα-Μοντέλο 7}
\begin{itemize}
    \item Σε αυτό το παράδειγμα στο κρυφό επίπεδο κάνουμε χρήση 6 επιπέδων Conv2D (6 συνελικτικά επίπεδα), το καθένα να έχει ένα Βatch Νormalization επίπεδο και 3 MaxPooling επίπεδα (ένα ανα 2 Conv2D επίπεδα). 
    Η έξοδος αυτών βρίσκει ενα Fully Connected Layer 256 νευρώνων
    \item Πλέον χρησιμοποιούμε ReduceLROnPlateau, Dropout μετά απο κάθε μπλοκ με διάφορες τιμές συν το Dropout που υπήρχε πριν (σύνολο δηλαδή 4 Dropout) και Weight Decay-L2 Regularization με l2(1e-4).
    \item Έχουμε 100 εποχές , 64 batch size και Adam Optimizer. \\ 
\end{itemize}
Ο κώδικας μας τυπώνει:
\begin{itemize}
    \item "Total training time = 553.50 seconds"
    \item "Test accuracy: 0.8631 , Test loss: 0.6295" \\ 
\end{itemize}
και έχουμε τα εξής διαγράμματα:

\begin{figure}[h]  
    \centering
    \includegraphics[width=0.5\textwidth]{7accuracy.png}
    \caption{Διάγραμμα ακρίβειας 7ου παραδείγματος}
\end{figure}
\newpage
\begin{figure}[h]  
    \centering
    \includegraphics[width=0.5\textwidth]{7loss.png}
    \caption{Διάγραμμα σφάλματος 7ου παραδείγματος}
\end{figure}
Σε αυτό το παράδειγμα έχουμε λοιπόν μια υγιή διαδικασία εκπαίδευσης σε σχέση με τα προηγούμενα
πολύπλοκα μοντέλα μας. Οι καμπύλες της ακρίβειας στην αρχή πηγαίνουν στις πρωτες 30 εποχές
εως και στο 0.8 μαζί (είναι δηλαδή η μία πολύ κοντά στην άλλη πράγμα που δηλώνει οτι το μοντέλο μαθαίνει
πολύ αποτελσματικά σε νέα δεδομένα). Έπειτα απο εκεί και πέρα η καμπύλη του training σταθεροποιείται
κάπου στο 0.9 και η validation στο 0.85 και αυτή  η μικρή διαφορά θεωρείται φυσιολογική και αποδεκτή. Ωστόσο το μοντέλο 
μετά απο λίγο παύει να βελτιώνεται και γύρω στις 85 εποχές πραγματοποιείται Early Stopping. \\ \\
Η ακρίβεια που πήραμε είναι η μέγιστη από όλα τα παραδείγματα που φτάσαμε και αυτό οφείλεται κυρίως 
στις πολλές βοηθητικές τεχνικές κανονικοποίησης που χρησιμοποιήσαμε. \\ \\
Στα επόμενα 2 παραδείγματα θα χρησιμοποιήσουμε μόνο ένα learning rate το οποίο όμως θα είναι αρκετά μικρότερο απο
το 0.001 που ξεκινάμε στα προηγούμενα παραδείγματα.

\section*{Παράδειγμα-Μοντέλο 8}
\begin{itemize}
    \item Σε αυτό το παράδειγμα στο κρυφό επίπεδο κάνουμε χρήση 6 επιπέδων Conv2D (6 συνελικτικά επίπεδα), το καθένα να έχει ένα Βatch Νormalization επίπεδο και 3 MaxPooling επίπεδα (ένα ανα 2 Conv2D επίπεδα). 
    Η έξοδος αυτών βρίσκει ενα Fully Connected Layer 256 νευρώνων
    \item Χρησιμοποιούμε  ένα Dropout(75\%) μόνο πάλι, Weight Decay-L2 Regularization με l2(1e-4) και learning rate ίσο με 5e-5 αυτή τη φορά.
    \item Έχουμε 100 εποχές , 64 batch size και Adam Optimizer. \\ 
\end{itemize}
Ο κώδικας μας τυπώνει:
\begin{itemize}
    \item "Total training time = 369.74 seconds"
    \item "Test accuracy: 0.7253 , Test loss: 0.9863" \\ 
\end{itemize}
και έχουμε τα εξής διαγράμματα:
\newpage

\begin{figure}[h]  
    \centering
    \includegraphics[width=0.5\textwidth]{8accuracy.png}
    \caption{Διάγραμμα ακρίβειας 8ου παραδείγματος}
\end{figure}
\begin{figure}[h]  
    \centering
    \includegraphics[width=0.5\textwidth]{8loss.png}
    \caption{Διάγραμμα σφάλματος 8ου παραδείγματος}
\end{figure}
Βλέπουμε πως λόγο του μικρού ρυθμού εκμάθησης που ρυθμίσαμε η καμπύλη της επικύρωσης στην ακρίβεια είναι πιο πάνω στην αρχή απο αυτή του 
training. Επίσης η χρήση του Dropout μειώνει το σκορ εκπαίδευσης και για αυτό είναι πιο κάτω. Μετά η καμπύλη training περνάει απο πάνω
και φτάνει ως και το 0.9 ενώ η άλλη μας καμπύλη βλέπουμε σταθεροποιείται γύρω στο 0.7. Έχουμε δηλαδή πάλι overfit. \\ \\
Η ακρίβεια λογικό να είναι πιο κάτω αφού ο ρυθμός εκμάθησης τώρα είναι σχετικά χαμηλός.


\section*{Παράδειγμα-Μοντέλο 9}
\begin{itemize}
    \item Σε αυτό το παράδειγμα στο κρυφό επίπεδο κάνουμε χρήση 6 επιπέδων Conv2D (6 συνελικτικά επίπεδα), το καθένα να έχει ένα Βatch Νormalization επίπεδο και 3 MaxPooling επίπεδα (ένα ανα 2 Conv2D επίπεδα). 
    Η έξοδος αυτών βρίσκει ενα Fully Connected Layer 256 νευρώνων
    \item Χρησιμοποιούμε  ένα Dropout(75\%) μόνο πάλι, Weight Decay-L2 Regularization με l2(1e-3) και learning rate ίσο με 1e-5 αυτή τη φορά.
    \item Έχουμε 150 εποχές , 32 batch size και Adam Optimizer. \\ 
\end{itemize}
Ο κώδικας μας τυπώνει:
\begin{itemize}
    \item "Total training time = 833.35 seconds"
    \item "Test accuracy: 0.6884 , Test loss: 1.5575" \\ 
\end{itemize}
και έχουμε τα εξής διαγράμματα:
\begin{figure}[h]  
    \centering
    \includegraphics[width=0.5\textwidth]{9accuracy.png}
    \caption{Διάγραμμα ακρίβειας 9ου παραδείγματος}
\end{figure}
\begin{figure}[h]  
    \centering
    \includegraphics[width=0.5\textwidth]{9loss.png}
    \caption{Διάγραμμα σφάλματος 9ου παραδείγματος}
\end{figure}
\newpage

Σε αντίθεση με όλα τα προηγούμενα παραδείγματα εδώ παρατηρούμε το φαινόμενο του underfit. Αυτό παρατηρείται στο διάγραμμα ακρίβειας από το 
γεγονός ότι η καμπύλη της επικύρωσης είναι για 60 εποχές πάνω απο την καμπύλη του training. Έπειτα η καμπύλη του training ξεπερνάει αυτή της επικύρωσης 
και συνεχίζει την ανοδική σταδιακή πορεία ενώ η άλλη σταθεροποιείται κάπου στο 0.65 χωρίς σημάδια βελτίωσης. Μετά εκεί στο τέλος μπορούμε να θεωρήσουμε
ότι έχουμε μια ήπια μορφή overfit. Η ακρίβεια λογικό να είναι τόσο χαμηλή αφού μειώσαμε το learning rate σε υπερβολικό βαθμό (1e-5).

\section*{Παραδείγματα Κατηγοριοποίησης}
Παρακάτω έχουμε τρία παραδείγματα κατηγοριοποίησης , τα δύο πρώτα αφορούν ένα αδύναμο μοντέλο, στα οποία στην αρχή προβλέπει σωστά και μετά όχι.
Στο τρίτο ένα καλύτερο μοντέλο πιάνει την ίδια εικόνα και βλέπουμε να προβλέπει σωστά τώρα σε σχέση με πριν. Ο κώδικας κάνει load το μοντέλο που θέλουμε
να δοκιμάσουμε , μετά εκτυπώνουμε την εικόνα που δοκιμάζουμε και ύστερα μέσω των συναρτήσεων της tensorflow βγάζουμε το αποτέλεσμα και το τυπώνουμε.

\begin{figure}[h]
    \centering
    \begin{minipage}{0.48\textwidth}
        \centering
        \includegraphics[width=0.75\linewidth]{classify1.png}
        \caption{Ορθή κατηγοριοποίηση/2ο μοντέλο}
    \end{minipage}\hfill 
    \begin{minipage}{0.48\textwidth}
        \centering
        \includegraphics[width=0.75\linewidth]{classify2.png}
        \caption{Λανθασμένη κατηγοριοποίηση/2ο μοντέλο}
    \end{minipage}
    \begin{minipage}{0.48\textwidth}
        \centering
        \includegraphics[width=0.75\linewidth]{classify3.png}
        \caption{Ορθή κατηγοριοποίηση/7ο μοντέλο}
    \end{minipage}
\end{figure}

\newpage
\section*{Τελικά Σχόλια για την Ολική Εργασία}
\begin{itemize}
    \item Παρατηρούμε οτι τα MLP μοντέλα μας, παρουσιάζουν σχετικά ομαλή εκπαίδευση ωστόσο μικρή ακρίβεια.
    \item Όσο πιο απλό είναι το μοντέλο μας τόσο μικρότερη θα είναι η ακρίβεια.
    \item Για αυτό και είναι εφικτό να δούμε ακρίβεια πιο μικρή και από τους αλγορίθμους ΚΝΝ και NC που χρησιμοποιήσαμε στην αρχή.
    \item Γενικά όμως η υλοποίηση ενός νευρωνικού δικτύου θεωρείται καλύτερη ειδικά όταν έχουμε πολύπλοκα δεδομένα όπως την Cifar10.
    \item Τα CNN μοντέλα μας αν και πετυχαίνουν μεγαλύτερη ακρίβεια (ειδικά όσο πιο πολύπλοκα γίνονται) παρουσιάζουν αρκετά το φαινόμενο του overfit.
    \item Καταλήξαμε ότι το όριο των συνελικτικών επιπέδων ήταν 6 (ή 3 μπλοκ) καθώς δεν παρατηρούσαμε παραπάνω βελτίωση.
    \item Ο σκοπός μετά ήταν με αυτά τα επίπεδα να λύσουμε αρκετά το φαινόμενο του overfit κρατώντας την ακρίβεια ψηλά πράγμα που καταφέραμε με τις διαφορές τεχνικές κανονικοποίησης.
    \item Το καλύτερο μας μοντέλο κατάφερε να κάνει υγιή εκπαίδευση και να φτάσει και 0.86 ακρίβεια.
    \item Τέλος είδαμε πως η χρήση πιο χαμηλού learning rate ενώ διώχνει το overfit μας ρίχνει χαμηλά και την ακρίβεια.
\end{itemize}
\end{document}

